\begingroup
\renewcommand{\arraystretch}{1.45}
\fancyhead[L]{\small BESTIMMUNG DER ANZIEHUNG}
\Needspace{42\baselineskip}
\AttractionArticle{14}

Durch Verbindung der Gleichungen 20, 21, 22 mit den Gleichungen des Art. 5 erhalten wir:
\[
\begin{aligned}
&at(A-a\cos E)=\alpha G-\alpha'G'\cos T-\alpha''G''\sin T\\[8pt]
&bt(B-b\sin E)=\beta G-\beta'G'\cos T-\beta''G''\sin T
\end{aligned}
\]
\par Setzt man daher zur Abkürzung
\[
\begin{aligned}
&(\alpha G-\alpha'G'\cos T-\alpha''G''\sin T)\\[8pt]
&\quad(\gamma-e\alpha+(\gamma'-e\alpha')\cos T+(\gamma''-e\alpha'')\sin T)=aX\\[8pt]
&(\beta G-\beta'G'\cos T-\beta''G''\sin T)\\[8pt]
&\quad(\gamma-e\alpha+(\gamma'-e\alpha')\cos T+(\gamma''-e\alpha'')\sin T)=bY\\[8pt]
&C(\gamma+\gamma'\cos T+\gamma''\sin T)\\[8pt]
&\quad(\gamma-e\alpha+(\gamma'-e\alpha')\cos T+(\gamma''-e\alpha'')\sin T)=Z
\end{aligned}
\]
so ergibt sich
\[
\begin{aligned}
&\mathrm d\xi=\frac{\varepsilon X\,\mathrm dT}{2\pi t^3\rho^3},\qquad
\mathrm d\eta=\frac{\varepsilon Y\,\mathrm dT}{2\pi t^3\rho^3},\qquad
\mathrm d\zeta=\frac{\varepsilon Z\,\mathrm dT}{2\pi t^3\rho^3}
\end{aligned}
\]
\Needspace{23\baselineskip}
\par Es gilt aber
\[
\begin{aligned}
&t\rho=\pm\surd{(G+G'\cos T^2+G''\sin T^2)}
\end{aligned}
\]
wobei das obere oder untere Vorzeichen gilt, je nachdem, ob $t$ eine positive oder eine negative Größe ist (denn $\rho$ wird seiner Natur nach stets positiv genommen), das heißt, je nachdem, ob der Koeffizient $\gamma$ positiv oder negativ ist. Hieraus folgt
\[
\begin{aligned}
&\frac{\varepsilon\,\mathrm dT}{2\pi t^3\rho^3}
=\pm\frac{\mathrm dT}{2\pi(G+G'\cos T^2+G''\sin T^2)^{\frac32}}
\end{aligned}
\]
wobei die Wahl des Vorzeichens vom Vorzeichen der Größe $\gamma\varepsilon$ abhängt.

\Needspace{27\baselineskip}
Um nun die Werte von $\xi,\eta,\zeta$ zu erhalten, müssen wir die Differentiale von dem Wert von $T$, dem $E=0$ entspricht, bis zu dem Wert integrieren, dem $E=360^0$ entspricht; oder auch, was offenbar auf dasselbe hinausläuft, von dem Wert von $T$, dem ein beliebiger Wert von $E$ entspricht, bis zu dem Wert, dem der Wert von $E$, vermehrt um $360^0$, entspricht. Wir dürfen daher von $T=0$ bis $T=360^0$ integrieren, sooft $\varepsilon\gamma$ eine positive Größe ist, oder von $T=360^0$ bis $T=0$, sooft $\varepsilon\gamma$ negativ ist. Somit gilt offenbar unabhängig vom Vorzeichen von $\varepsilon\gamma$:
\[
\begin{aligned}
&\xi=\int\frac{X\,\mathrm dT}{2\pi(G+G'\cos T^2+G''\sin T^2)^{\frac32}}\\[8pt]
&\eta=\int\frac{Y\,\mathrm dT}{2\pi(G+G'\cos T^2+G''\sin T^2)^{\frac32}}\\[8pt]
&\zeta=\int\frac{Z\,\mathrm dT}{2\pi(G+G'\cos T^2+G''\sin T^2)^{\frac32}}
\end{aligned}
\]
wobei die Integrationen von $T=0$ bis $T=360^0$ zu erstrecken sind.
\SourceLeaf{360}
\Needspace{43\baselineskip}
\AttractionArticle{15}

Man erkennt ohne Schwierigkeit, dass die Integrale
\[
\begin{aligned}
&\int\frac{\cos T\,\mathrm dT}{(G+G'\cos T^2+G''\sin T^2)^{\frac32}}\\[8pt]
&\int\frac{\sin T\,\mathrm dT}{(G+G'\cos T^2+G''\sin T^2)^{\frac32}}\\[8pt]
&\int\frac{\cos T\sin T\,\mathrm dT}{(G+G'\cos T^2+G''\sin T^2)^{\frac32}}
\end{aligned}
\]
von $T=180^0$ bis $T=360^0$ erstreckt Werte erhalten, die denen gleich sind, die sie bei Erstreckung von $T=0$ bis $T=180^0$ erhalten, jedoch mit entgegengesetzten Vorzeichen; daher werden diese Integrale, von $T=0$ bis $T=360^0$ erstreckt, offenbar $=0$. Hieraus schließen wir. dass
\[
\begin{aligned}
&\xi{}=\int\frac{\begin{gathered}((\gamma-e\alpha)\alpha G-(\gamma'-e\alpha')\alpha'G'\cos T^2\\-(\gamma''-e\alpha'')\alpha''G''\sin T^2)\,\mathrm dT\end{gathered}}{2\pi a(G+G'\cos T^2+G''\sin T^2)^{\frac32}}\\[8pt]
&\eta{}=\int\frac{\begin{gathered}((\gamma-e\alpha)\beta G-(\gamma'-e\alpha')\beta'G'\cos T^2\\-(\gamma''-e\alpha'')\beta''G''\sin T^2)\,\mathrm dT\end{gathered}}{2\pi b(G+G'\cos T^2+G''\sin T^2)^{\frac32}}\\[8pt]
&\zeta{}=\int\frac{\begin{gathered}((\gamma-e\alpha)\gamma+(\gamma'-e\alpha')\gamma'\cos T^2\\+(\gamma''-e\alpha'')\gamma''\sin T^2)C\,\mathrm dT\end{gathered}}{2\pi(G+G'\cos T^2+G''\sin T^2)^{\frac32}}
\end{aligned}
\]
\Needspace{38\baselineskip}
wobei die Integrale von $T=0$ bis $T=360^0$ zu erstrecken sind. Bezeichnen wir daher die Werte der über dasselbe Intervall erstreckten Integrale
\[
\begin{aligned}
&\int\frac{\cos T^2\,\mathrm dT}{2\pi((G+G')\cos T^2+(G+G'')\sin T^2)^{\frac32}}\\[8pt]
&\int\frac{\sin T^2\,\mathrm dT}{2\pi((G+G')\cos T^2+(G+G'')\sin T^2)^{\frac32}}
\end{aligned}
\]
mit $P,Q$, so ergibt sich
\[
\begin{aligned}
&a\xi=((\gamma-e\alpha)\alpha G-(\gamma'-e\alpha')\alpha'G')P\\[6pt]
&\qquad{}+((\gamma-e\alpha)\alpha G-(\gamma''-e\alpha'')\alpha''G'')Q\\[8pt]
&b\eta=((\gamma-e\alpha)\beta G-(\gamma'-e\alpha')\beta'G')P\\[6pt]
&\qquad{}+((\gamma-e\alpha)\beta G-(\gamma''-e\alpha'')\beta''G'')Q\\[8pt]
&\zeta=((\gamma-e\alpha)\gamma+(\gamma'-e\alpha')\gamma')CP\\[6pt]
&\qquad{}+((\gamma-e\alpha)\gamma+(\gamma''-e\alpha'')\gamma'')CQ
\end{aligned}
\]
womit unser Problem vollständig gelöst ist.
\SourceLeaf{361}
\clearpage
\endgroup
